A implementação dos algoritmos genéticos no sistema IGRIS representa a aplicação prática dos conceitos apresentados nos capítulos anteriores. Enquanto a fundamentação teórica discutiu os princípios dos Algoritmos Genéticos e sua utilização em problemas de busca e otimização, este capítulo apresenta como esses conceitos foram empregados, de maneira concreta, na geração automatizada de grades horárias acadêmicas.

O problema tratado pelo sistema está inserido no contexto do \textit{University Course Timetabling Problem} (UCTP), que envolve a distribuição de demandas acadêmicas em horários disponíveis, considerando restrições como turma, turno, professor habilitado, disponibilidade docente e conflitos de alocação \cite{premananda2024}. Trata-se de um problema combinatório, pois o número de combinações possíveis cresce significativamente conforme aumentam a quantidade de turmas, disciplinas, professores e horários, característica recorrente em problemas de timetabling educacional \cite{ceschia2023}. Diante desse cenário, uma busca exaustiva torna-se pouco viável, o que justifica o uso de métodos aproximativos e meta-heurísticos, como os Algoritmos Genéticos \cite{eiben2015}.

No IGRIS, essa abordagem foi implementada por meio de diferentes serviços de otimização, denominados Sunny, Morpheus e Rocket, responsáveis pela geração automatizada de grades horárias a partir dos dados acadêmicos enviados pelo núcleo do sistema. Embora todos compartilhem o mesmo objetivo de construir, avaliar e retornar a melhor grade encontrada com base em uma função \textit{fitness} e em operadores evolutivos, cada serviço organiza internamente esse processo de maneira distinta, com diferenças na representação dos dados, na formação da população inicial, na avaliação das soluções e no tratamento de conflitos.

Como discutido anteriormente, Algoritmos Genéticos não têm como finalidade garantir, necessariamente, a obtenção da solução ótima global. Seu objetivo é encontrar soluções satisfatórias dentro de um tempo computacional aceitável, especialmente em espaços de busca extensos e complexos \cite{goldberg1989}. Essa característica está presente nos algoritmos do IGRIS, que trabalham com sucessivas gerações de grades candidatas, avaliando sua qualidade, preservando boas soluções e explorando novas combinações por meio de seleção, recombinação e mutação. Assim, a resposta final corresponde à melhor grade identificada durante a execução, considerando as restrições acadêmicas e os parâmetros configurados no sistema.